\chapter{VFS：super\_block、inode 与 dentry}
\label{chap:vfs}
T113 版 miniOS 的 VFS 已经从“挂载表 + 整条路径交给后端”改成 Linux 式的对象模型：每次挂载生成一个 \code{super\_block}，路径按分量逐级建立 \code{dentry}，每个\code{dentry} 指向一个 \code{inode}，打开的文件是 \code{struct file}，挂载点之间用\code{vfs\_mount} 连成树。本章按当前代码说明这几个对象怎么来、谁持有谁、一次\code{open/read/write} 怎样穿过它们，以及和 Linux、xv6 的差别。

源码：

\begin{table}[H]
\centering
\small
\begin{tabularx}{\linewidth}{>{\raggedright\arraybackslash}X>{\raggedright\arraybackslash}X}
\toprule
文件 & 内容 \\
\midrule
\code{src/include/fs/vfs.h} & 对象定义、三张操作表、后端接口 \code{vfs\_backend\_ops} \\
\code{src/kernel/fs/vfs.c} & 对象池、\code{lookup\_path()}、挂载、open/read/write/close/readdir \\
\code{src/kernel/fs/minifs.c} & \file{/} 根文件系统（SD 卡分区 2，只读） \\
\code{src/kernel/fs/fat32.c} & \file{/boot}（SD 卡分区 1，读写，带 create/truncate/close） \\
\code{src/kernel/fs/bootfs.c} & \file{/boot} 的备用后端（SPL 预加载镜像，只读） \\
\code{src/kernel/fs/devfs.c} + \code{chrdev.c} & \file{/dev}，名字 → 设备号 → 驱动 \\
\code{src/kernel/syscall.c} & \code{open/read/write/close/dup/readdir/exec} 进入 VFS \\
\bottomrule
\end{tabularx}
\end{table}

\section{改了什么}

\begin{table}[H]
\centering
\small
\begin{tabularx}{\linewidth}{l>{\raggedright\arraybackslash}X>{\raggedright\arraybackslash}X}
\toprule
方面 & 旧实现（d8a6613） & 现在 \\
\midrule
挂载 & \code{mounts[4]}，每项只有路径和 \code{ops} & \code{vfs\_mount} + \code{super\_block}，子挂载记下父挂载和挂载点 dentry \\
名字解析 & \code{ops->lookup("/boot/fw\_xr829.bin")} 一次解析整条路径 & \code{lookup\_path()} 按分量走，每一级一个 \code{dentry} \\
文件身份 & \code{open\_files[]} 里复制一份 \code{vfs\_node} & 共享的 \code{inode}，同一文件多次打开指向同一个 inode \\
操作分派 & 直接调后端函数 & \code{inode->i\_op} / \code{file->f\_op}，再经适配层到后端 \\
文件位置 & 全局 \code{vfs\_lock} 下更新 offset & 每个 \code{file} 一把 \code{f\_pos\_lock}，字符设备不加锁 \\
控制台 & fd 0/1/2 在系统调用里特判 UART & \file{/init} 打开 \file{/dev/console} 再 \code{dup}，全走 VFS \\
\bottomrule
\end{tabularx}
\end{table}

后端接口 \code{struct vfs\_backend\_ops} 还保留着，每个文件系统仍然接收\textbf{累计的绝对路径}。新对象层和旧后端之间由 \code{vfs.c} 里的一组通用 \code{backend\_*} 函数衔接，后面第~\ref{sec:vfs-ops} 节专门讲。

\section{五种对象和它们的池}

所有对象都是 \code{vfs.c} 里的静态数组，没有动态分配。\code{used} 表示槽位占用，\code{references} 是引用计数。

\begin{table}[H]
\centering
\small
\begin{tabularx}{\linewidth}{>{\raggedright\arraybackslash}X>{\raggedright\arraybackslash}Xl>{\raggedright\arraybackslash}X}
\toprule
对象 & 池 & 查找键 & 谁持有引用 \\
\midrule
\code{struct vfs\_mount} & \code{mounts[8]} & 挂载路径字符串 & 永不释放（没有 umount） \\
\code{struct super\_block} & \code{super\_blocks[8]} & — & 挂载 1 个，加上它名下每个 inode 各 1 个 \\
\code{struct dentry} & \code{dentry\_cache[128]} & \code{(d\_sb, d\_path)} & 子 dentry、\code{file}、\code{vfs\_mount}、正在进行的 lookup \\
\code{struct inode} & \code{inode\_cache[128]} & \code{(i\_sb, i\_ino)} & 指向它的 dentry \\
\code{struct file} & \code{open\_files[32]} & — & 进程的 fd 表（\code{dup}、\code{fork} 会加）和进行中的 read/write \\
\bottomrule
\end{tabularx}
\end{table}

全部对象都由一把全局 \code{struct mutex vfs\_lock} 保护。真正的 I/O 不在这把锁里做（见第~\ref{sec:vfs-io} 节）。

\subsection{字段}

\begin{ccode}
struct super_block {               /* 一次挂载的文件系统实例 */
    uint32_t s_magic;              /* minifs 0x10203040, fat32 0x4d44,
                                      devfs 0x1373, bootfs 0x424f4f54 */
    const struct super_operations *s_op;
    struct dentry *s_root;         /* 本章件系统的根 dentry */
    void *s_fs_info;               /* minifs &fs_sb, fat32 &fs, devfs nodes */
    const struct vfs_backend_ops *s_backend;   /* 旧后端接口 */
};

struct inode {                     /* 一个文件或目录本身 */
    uint32_t i_ino, i_size, i_type;           /* DIR=1 FILE=2 CHAR=3 */
    struct super_block *i_sb;
    const struct inode_operations *i_op;      /* lookup/create/truncate */
    const struct file_operations *i_fop;      /* 打开后用的操作表 */
    void *i_private;               /* 后端私有（fat32: fat_inode 槽位） */
    struct vfs_node i_backend_node;           /* 交给后端的句柄 */
};

struct dentry {                    /* 路径中的一个名字 */
    char d_name[31];               /* 本级名字，如 "fw_xr829.bin" */
    char d_path[128];              /* 累计绝对路径，如 "/boot/fw_xr829.bin" */
    struct dentry *d_parent;       /* 持有父 dentry 的一个引用 */
    struct inode *d_inode;         /* 持有 inode 的一个引用 */
    struct super_block *d_sb;
};

struct file {                      /* 一次 open 的结果 */
    uint32_t f_flags, f_pos;
    struct dentry *f_dentry;       /* 持有 dentry 的一个引用 */
    struct inode *f_inode;
    const struct file_operations *f_op;
    struct mutex f_pos_lock;
};

struct vfs_mount {                 /* 挂载树的一条边 */
    char path[128];
    struct super_block *sb;
    struct dentry *mnt_root;       /* == sb->s_root */
    struct vfs_mount *parent;      /* 被挂在哪个挂载下面 */
    struct dentry *mountpoint;     /* 父文件系统里的那个目录 dentry */
};
\end{ccode}

\section{启动后的对象图}

\code{kernel\_boot()} 依次挂 \file{/}（minifs）、\file{/boot}（FAT32，失败退回 bootfs）、\file{/dev}（devfs）。\file{/init} 打开 \file{/dev/console} 并 \code{dup} 两次之后，对象之间是这样连的（括号里是引用计数；这是设计上的结果，前提是第~\ref{sec:vfs-issues} 节第 1 条的minifs 目录查找问题已经修好）：

\begin{txtcode}
mounts[0] "/"          mounts[1] "/boot"                mounts[2] "/dev"
  | parent=0             | parent=mounts[0]               | parent=mounts[0]
  |                      | mountpoint ---------+          | mountpoint -------------+
  v                      v                     |          v                         |
sb0 minifs (4)         sb1 fat32 (2)           |        sb2 devfs (3)               |
  | s_root               | s_root              |          | s_root                  |
  v                      v                     |          v                         |
dentry "/" sb0 (3) <-+  dentry "/boot" sb1 (1) |        dentry "/dev" sb2 (2) <-+   |
  | d_inode          |    | d_inode            |          | d_inode             |   |
  v                  |    v                    |          v                     |   |
inode ino=0 DIR      |  inode ino=0 DIR        |        inode ino=0 DIR         |   |
                     |                         |                                |   |
     d_parent -------+                         |      d_parent -----------------+   |
                     |                         |                                |   |
dentry "/boot" sb0 (1) <-----------------------+      dentry "/dev/console" sb2 (1)  |
  +> inode ino=4 DIR (minifs)                           +> inode ino=0x10000 CHAR    |
dentry "/dev" sb0 (1) <--------------------------------------------------------------+
  +> inode ino=3 DIR (minifs)                                     ^ f_dentry
                                                                  |
                                    open_files[k] (3) ------------+
                                    ^     ^     ^
                         init: fd 0  fd 1  fd 2         (fork 后子进程再各加 1)
\end{txtcode}

几个要点：

\begin{itemize}
  \item 同一个路径字符串可以对应两个 dentry：\code{"/boot"} 在 sb0（minifs 里的空目录，也就是挂载点）和 sb1（FAT32 的根）各有一个。dentry 的键是\code{(d\_sb, d\_path)}，所以两者不会撞。
  \item 挂载根的 inode 号统一是 0，类型是目录，不是后端里的真实根 inode（minifs 的根是 1）。后端靠路径而不是 inode 号处理根目录，所以没影响。
  \item sb0 的计数 4 = 挂载本身 1 + 根 inode 1 + 两个挂载点 inode 各 1。sb0 根dentry 的计数 3 = 挂载 1 + 两个挂载点子 dentry 各 1。
  \item 设备文件的 inode 号就是设备号 \code{MINIOS\_MKDEV(1, 0) = 0x10000}。
\end{itemize}

\section{路径查找：\code{lookup\_path()}}

\begin{ccode}
static struct dentry *lookup_path(const char *path);
\end{ccode}

返回的 dentry 带一个调用者的引用，用完要 \code{dput()}。过程：

\begin{enumerate}
  \item \code{find\_mount(path)}：在 \code{mounts[]} 里取最长的匹配前缀，并要求前缀后面是\file{/} 或字符串结束，所以 \file{/bootx} 只匹配 \file{/}。
  \item 从这个挂载的 \code{mnt\_root} 出发，引用加一。
  \item 跳过挂载路径，剩下的部分逐个分量处理：- 名字长度为 0 或超过 \code{VFS\_NAME\_MAX}（30）就失败；- \code{append\_component()} 把分量接到累计路径 \code{absolute} 上；- \code{alloc\_dentry(sb, current, name, absolute)}：先按 \code{(sb, absolute)} 找已有的 dentry，找到就加引用；没有就占一个空槽，\code{d\_parent = current} 并给父dentry 加引用；- 新 dentry 还没有 inode 时，调用 \code{current->d\_inode->i\_op->lookup(dir, child)}，失败就把这个 dentry 放掉，整次查找失败；- \code{dput(current)}，\code{current = child}。
  \item 返回最后一级。
\end{enumerate}

以 \code{t113\_vfs\_open("/boot/fw\_xr829.bin")} 为例：

\begin{txtcode}
find_mount("/boot/fw_xr829.bin")             "/" 和 "/boot" 都匹配，取 "/boot"
current = mounts[1].mnt_root                 dentry "/boot" (sb1)
absolute = "/boot"
分量 "fw_xr829.bin"
+-- absolute = "/boot/fw_xr829.bin"
+-- alloc_dentry(sb1, "/boot", ...)          dcache 未命中，新建，父 dentry 引用 +1
+-- i_op->lookup = backend_lookup()
|   +-- fat32_ops.lookup("/boot/fw_xr829.bin", &node)
|       +-- resolve() → iget()               fat_inode 槽位 ref +1
|           node = { id = 槽位地址, size, FILE, private_data = 槽位 }
|   +-- alloc_inode(sb1, &node)              按 (sb1, id) 查 icache，未命中则新建
+-- dput("/boot")
返回 dentry "/boot/fw_xr829.bin"
\end{txtcode}

\file{/bin/sh} 这类根文件系统里的多级路径会走两步：先为 \code{"/bin"} 建 dentry，由 minifs 查出目录 inode，再为 \code{"/bin/sh"} 建 dentry。（minifs 当前在第一步会失败，见第~\ref{sec:vfs-issues} 节。）

\textbf{跨挂载点的方式。} 查找从不“走过”挂载点：选哪个文件系统完全由第 1 步的字符串前缀决定，挂载点 dentry 只在挂载时用来确认父文件系统里确实有这个目录。Linux 是在逐级走的过程中，看到 dentry 被挂载过（\code{d\_mounted}）就切到被挂载文件系统的根（\code{traverse\_mounts()}），\code{..} 走到挂载根时再沿\code{mnt\_parent} 回到挂载点。miniOS 现在不处理 \code{.} 和 \code{..}：它们作为普通分量交给后端，\file{/boot/..} 回不到 minifs 的 \file{/}。

\section{三张操作表和适配层}\label{sec:vfs-ops}

\code{vfs.h} 定义了 Linux 形状的三张表：

\begin{ccode}
struct super_operations { void (*put_super)(struct super_block *); };
struct inode_operations {
    int (*lookup)(struct inode *dir, struct dentry *child);
    int (*create)(struct inode *dir, struct dentry *child);
    int (*truncate)(struct inode *inode);
};
struct file_operations {
    int  (*open)(struct inode *, struct file *);
    int  (*read)(struct file *, void *, uint32_t);
    int  (*write)(struct file *, const void *, uint32_t);
    int  (*readdir)(struct inode *, const char *path, uint32_t index,
                    struct vfs_dirent *);
    void (*release)(struct inode *, struct file *);
};
\end{ccode}

各文件系统还没有自己的这三张表。\code{vfs.c} 给每个 inode 装的都是同一组\code{backend\_iops / backend\_fops / backend\_sops}，它们再通过 \code{i\_sb->s\_backend}调用旧接口：

\begin{txtcode}[basicstyle=\ttfamily\scriptsize]
VFS 对象层                         适配层（vfs.c）              后端（每个文件系统）
inode->i_op->lookup(dir, child) -> backend_lookup()  ---------> ops->lookup(child->d_path, &node)
inode->i_op->create(dir, child) -> backend_create()  ---------> ops->create(child->d_path, &node)
inode->i_op->truncate(inode)    -> backend_truncate() --------> ops->truncate(&i_backend_node)
file->f_op->open(inode, file)   -> backend_open()    ---------> ops->open(&i_backend_node, flags)
file->f_op->read(file, buf, n)  -> backend_read()    ---------> ops->read(&i_backend_node, f_pos, buf, n)
file->f_op->write(file, buf, n) -> backend_write()   ---------> ops->write(&i_backend_node, f_pos, buf, n)
inode->i_fop->readdir(...)      -> backend_readdir() ---------> ops->readdir(path, index, entry)
iput() 最后一个引用              -----------------------------> ops->close(&i_backend_node)
\end{txtcode}

\code{lookup} 的参数已经是 Linux 的 \code{(父 inode, 子 dentry)}，但适配层只用了\code{child->d\_path}，后端还是从文件系统的根把整条路径重新解析一遍。把某个文件系统改成真正的逐级查找，只需要给它写自己的 \code{inode\_operations.lookup}：用\code{dir->i\_ino} 和 \code{child->d\_name} 在一个目录里找一项，然后在 \code{alloc\_inode} 时装上它自己的表。

\subsection{各后端填的 \code{vfs\_node}}

\begin{table}[H]
\centering
\small
\begin{tabularx}{\linewidth}{l>{\raggedright\arraybackslash}Xl>{\raggedright\arraybackslash}X}
\toprule
后端 & \code{id}（成为 \code{i\_ino}） & \code{private\_data} & 实现的可选操作 \\
\midrule
minifs & 磁盘 inode 号 & 0 & 只有 lookup/read/readdir \\
fat32 & \code{fat\_inode} 槽位地址 & 同一槽位 & create、truncate、write、close \\
bootfs & 镜像在 DDR 中的地址 & 0 & 只有 lookup/read \\
devfs & \code{MINIOS\_MKDEV(major, minor)} & \code{\&nodes[i]} & open、write \\
\bottomrule
\end{tabularx}
\end{table}

FAT32 用槽位地址当 inode 号，因为空文件的首簇都是 0，不能当身份。\code{fat32\_lookup()} 每次都会对 \code{fat\_inode} 加一次引用；如果 VFS 的 icache 里已经有这个 inode，\code{alloc\_inode()} 立刻用 \code{ops->close()} 把这次多出来的引用还回去，保证“一个 VFS inode 对应 FAT32 的一个引用”。

\section{挂载：\code{t113\_vfs\_mount\_fs()}}

\begin{ccode}
int t113_vfs_mount_fs(const char *path, uint32_t magic, void *fs_info,
                      const struct vfs_backend_ops *ops);
\end{ccode}

旧的 \code{t113\_vfs\_mount(path, ops)} 仍在，等于 \code{magic = 0, fs\_info = 0}。步骤：

\begin{txtcode}
t113_vfs_mount_fs("/boot", 0x4d44, &fs, &fat32_ops)
+-- 检查 path 以 '/' 开头，ops->lookup 必须有
+-- lock vfs_lock
+-- 不是 "/" 时：
|   +-- parent_mount = find_mount("/boot")         → mounts[0] "/"
|   +-- mountpoint   = lookup_path("/boot")        → minifs 里的 dentry，必须是目录
+-- 路径已挂载过 → 失败
+-- alloc_super()                                   sb 引用 = 1（归挂载）
|   s_backend = ops, s_magic = magic, s_fs_info = fs_info
+-- alloc_inode(sb, {id 0, size 0, DIR})            根 inode，sb 引用 +1
+-- alloc_dentry(sb, NULL, "boot", "/boot")         根 dentry，没有父
+-- root->d_inode = 根 inode；sb->s_root = root
+-- mounts[i] = { path, sb, mnt_root = root, parent = parent_mount,
                  mountpoint = mountpoint }         挂载点 dentry 的引用一直留着
\end{txtcode}

所以挂载顺序有硬性要求：子挂载点必须是父文件系统里已经存在的目录。\code{tools/mkfs\_minifs.c} 因此在镜像里建了 \file{/bin}、\file{/dev}，现在又加了 \file{/boot}。\file{/dev} 的 devfs 也要求 minifs 里先有 \file{/dev} 目录。

\section{open / read / write / close}\label{sec:vfs-io}

\subsection{系统调用入口}

\begin{txtcode}
用户 open(path, flags)
+-- sys_open()                                      kernel/syscall.c
    +-- flags 只能是 RDONLY/RDWR，加 O_CREATE(0x200)、O_TRUNC(0x400)
    +-- copy_user_string(path)
    +-- t113_vfs_open(path, flags)
    +-- 放进 t113_current->files[] 第一个空位（16 个 fd）
\end{txtcode}

\code{read/write} 先检查 fd 和用户缓冲区（\code{t113\_user\_vm\_range\_writable}、\code{user\_range\_ok}），然后直接调 \code{t113\_vfs\_read/write}。fd 0/1/2 不再特判：\file{/init} 一开始 \code{open("/dev/console", O\_RDWR)} 得到 fd 0，再 \code{dup(0)} 两次得到1 和 2。\code{dup} 调 \code{t113\_vfs\_file\_get()}，\code{fork} 时 \code{t113\_kernel\_clone()} 复制 fd 表，也对每个 fd 调一次。

\subsection{\code{t113\_vfs\_open()}}

\begin{txtcode}
t113_vfs_open(path, flags)
+-- 参数检查：O_TRUNC 只能和 RDWR 一起用
+-- lock vfs_lock
+-- dentry = lookup_path(path)
+-- 找不到且有 O_CREATE → create_path(path)
|   +-- split_parent()：  "/boot/a.txt" → "/boot" + "a.txt"
|   +-- parent = lookup_path("/boot")，必须是目录
|   +-- child = alloc_dentry(sb, parent, "a.txt", path)
|   +-- parent->d_inode->i_op->create(dir, child)    → fat32_create()
+-- dentry 不存在或是目录 → 失败
+-- O_TRUNC → inode->i_op->truncate(inode)           → fat32_truncate()
+-- 在 open_files[] 占一个槽：
|   f_flags = mode, f_pos = 0, f_dentry = dentry（接过 lookup 的那个引用）
|   f_inode = inode, f_op = inode->i_fop, 初始化 f_pos_lock
+-- 以 RDWR 打开但后端没有 write → 失败（minifs、bootfs 都是只读）
+-- f_op->open(inode, file)                          devfs 只检查 flags
+-- unlock vfs_lock，返回 file
\end{txtcode}

\subsection{\code{t113\_vfs\_read()} / \code{t113\_vfs\_write()}}

\begin{txtcode}
t113_vfs_read(file, buf, n)
+-- lock vfs_lock
|   +-- 检查 file 有效、后端有 read
|   +-- file->references++                ← 临时引用
|   +-- positional = (i_type != CHAR)
+-- unlock vfs_lock
+-- positional ? lock f_pos_lock
+-- result = f_op->read(file, buf, n)     → backend_read → ops->read(node, f_pos, ...)
+-- result > 0 且 positional → f_pos += result
+-- positional ? unlock f_pos_lock
+-- t113_vfs_close(file)                  ← 放掉临时引用
\end{txtcode}

\code{write} 相同，另外在 \code{f\_pos} 超过 \code{i\_size} 时同时更新 \code{i\_size} 和\code{i\_backend\_node.size}。

这里有两处设计：

\begin{itemize}
  \item \textbf{I/O 不持有 \code{vfs\_lock}。} 控制台读会睡眠等按键，SD 读要等卡，如果在全局锁里做，一个进程读控制台就会卡住所有进程的 open。代价是 I/O 期间别的线程可能 \code{close} 这个 fd，所以读写前先拿一个临时引用，保证 \code{file} 槽位在 I/O结束前不被回收。
  \item \textbf{字符设备不加 \code{f\_pos\_lock}。} 设备没有位置，后端也忽略 offset；共享同一个控制台 \code{file} 的父子进程不应该因为一个在等输入，另一个就不能输出。Linux 的做法一样：只有带 \code{FMODE\_ATOMIC\_POS} 的普通文件才在 \code{fdget\_pos()} 里锁\code{f\_pos\_lock}。
\end{itemize}

\subsection{\code{t113\_vfs\_close()} 与释放链}

\begin{txtcode}
t113_vfs_close(file)
+-- --file->references == 0
    +-- f_op->release（适配层为空）
    +-- dput(f_dentry)
    |   +-- --references == 0
    |       +-- iput(d_inode)
    |       |   +-- --references == 0
    |       |       +-- i_private 非空且后端有 close → ops->close()
    |       |       |   （fat32：dirty 时回写目录项，fat_inode ref--）
    |       |       +-- sb->references--，槽位释放
    |       +-- 槽位释放
    |       +-- dput(d_parent)       一路向上，直到遇到还有别人引用的 dentry
    +-- 槽位释放
\end{txtcode}

dentry 的引用归零就立即释放，没有保留在 LRU 里，也没有负 dentry。所以“dcache”只在同一对象同时被多人使用时起作用：\file{/dev/console} 被三个 fd 共享时只有一个 dentry；关掉以后再次打开，会重新逐级调用后端 lookup。

\subsection{\code{t113\_vfs\_readdir()}}

\begin{txtcode}
t113_vfs_readdir(path, index, entry)
+-- lock vfs_lock
+-- dentry = lookup_path(path)，必须是目录
+-- i_fop->readdir(inode, path, index, entry)   → ops->readdir(path, index, entry)
+-- dput(dentry)
+-- unlock vfs_lock
\end{txtcode}

\code{ls} 每读一项就调一次，后端每次从头数到第 \code{index} 项。\code{readdir} 是在\code{vfs\_lock} 里做的，FAT32 读目录簇时会让其他 VFS 调用等待。

\section{引用计数规则一览}

\begin{table}[H]
\centering
\small
\begin{tabularx}{\linewidth}{>{\raggedright\arraybackslash}X>{\raggedright\arraybackslash}X>{\raggedright\arraybackslash}X}
\toprule
动作 & 加引用 & 减引用 \\
\midrule
\code{alloc\_super()} & sb = 1（挂载） & — \\
\code{alloc\_inode()} 新建 & inode = 1（dentry），sb +1 & — \\
\code{alloc\_inode()} 命中 & inode +1 & 后端多拿的那个引用立即 close \\
\code{alloc\_dentry()} 新建 & dentry = 1，父 dentry +1 & — \\
\code{alloc\_dentry()} 命中 & dentry +1 & — \\
\code{lookup\_path()} 每走一步 & child & \code{dput(current)} \\
\code{t113\_vfs\_open()} & file = 1，接管 dentry 的引用 & 失败时 \code{dput} \\
\code{dup} / \code{fork} & file +1 & — \\
read/write 期间 & file +1 & 结束时 −1 \\
\code{close} / 进程退出 & — & file −1 → \code{dput} → \code{iput} → 后端 close \\
\bottomrule
\end{tabularx}
\end{table}

\section{与 Linux、xv6 对照}

\begin{table}[H]
\centering
\small
\begin{tabularx}{\linewidth}{l>{\raggedright\arraybackslash}X>{\raggedright\arraybackslash}X>{\raggedright\arraybackslash}X}
\toprule
方面 & xv6 & Linux & T113 miniOS \\
\midrule
有没有 VFS & 没有，只有一种文件系统 & 有 & 有，四种后端 \\
超级块 & 磁盘结构 \code{struct superblock}，全局一份 & \code{struct super\_block}，\code{s\_op}、\code{s\_root}、\code{s\_fs\_info} & 同名结构，\code{s\_backend} 指向旧后端 \\
inode 缓存 & \code{itable.inode[NINODE=50]}，键 \code{(dev, inum)} & icache 哈希，键 \code{(sb, ino)} & \code{inode\_cache[128]} 线性扫描，键 \code{(sb, ino)} \\
名字缓存 & 无，\code{namex()} 每次都调 \code{dirlookup()} & dcache，键 \code{(父 dentry, 名字)}，LRU + 负 dentry，RCU 遍历 & \code{dentry\_cache[128]}，键 \code{(sb, 完整路径)}，引用归零即释放 \\
逐级查找 & \code{namex()} → \code{dirlookup(dp, name)} & \code{link\_path\_walk()} → \code{i\_op->lookup(dir, dentry)} & \code{lookup\_path()} 逐级建 dentry，后端仍收整条路径 \\
跨挂载点 & 无挂载 & \code{traverse\_mounts()}，\code{..} 经 \code{mnt\_parent} 回退 & 最长前缀选挂载；\code{..} 不跨挂载 \\
挂载 & 无 & \code{struct mount}：\code{mnt\_parent}、\code{mnt\_mountpoint}、\code{mnt\_root} & \code{vfs\_mount}：\code{parent}、\code{mountpoint}、\code{mnt\_root} \\
打开文件 & \code{ftable.file[100]}，\code{ref}、\code{off} & \code{struct file}，\code{f\_count}、\code{f\_pos}、\code{f\_pos\_lock} & \code{open\_files[32]}，\code{references}、\code{f\_pos}、\code{f\_pos\_lock} \\
设备文件 & 磁盘 \code{T\_DEVICE} inode + \code{devsw[major]} & 字符设备 inode 的 \code{i\_fop} 换成 \code{def\_chr\_fops} → \code{cdev} & devfs 节点 + \code{chrdev\_table[32]} 按设备号分派 \\
锁 & 每个 inode 一把 sleeplock & \code{i\_rwsem}、\code{d\_lock}、RCU 等细粒度锁 & 全局 \code{vfs\_lock} + 每个 file 的 \code{f\_pos\_lock} \\
\bottomrule
\end{tabularx}
\end{table}

\section{现状与已知问题}\label{sec:vfs-issues}

以下是读代码得出的结论，还没有上板验证。

\begin{enumerate}
  \item \textbf{minifs 的 lookup 不接受目录，启动会在挂 \file{/boot} 时停下。} \code{minifs\_vfs\_lookup()} 调的是 \code{t113\_minifs\_open()}，后者要求\code{inode.type == MINIFS\_T\_FILE}。旧 VFS 把整条路径交给后端，不需要单独查目录，所以没出问题。现在 \code{t113\_vfs\_mount\_fs("/boot")} 要先\code{lookup\_path("/boot")} 拿挂载点，minifs 返回失败，FAT32 和 bootfs 都挂不上，\code{kernel\_boot()} 会 panic \code{cannot mount boot filesystem at /boot}。即使绕过这一步，\file{/dev} 的挂载和 \code{exec("/bin/sh")} 也会因为 \file{/dev}、\file{/bin} 查不到而失败。改法：\code{minifs\_vfs\_lookup()} 直接用 \code{namei()} + \code{read\_inode()}，不检查类型。
  \item \textbf{后端仍按整条路径解析。} 每走一级，minifs 都从根 inode 重新 \code{namei()}，深度为 d 的路径要做大约 d²/2 次目录查找。改成真正的逐级 lookup 后，每级只查一个目录。
  \item \textbf{dentry 不缓存。} 引用归零就释放；改进方向是保留未使用的 dentry（LRU）并缓存“不存在”的结果。
  \item \textbf{名字长度上限 30。} \code{VFS\_NAME\_MAX} 取自 minifs 的 \code{DIRSIZ}。FAT32 支持127 字节的长文件名，但超过 30 字节的分量会在 \code{lookup\_path()} 里直接失败。\code{d\_path} 上限 128 字节也限制了路径总长。
  \item \textbf{不处理 \code{.}、\code{..} 和跨挂载回退}，也没有每进程 cwd（\code{getcwd} 返回\code{"/"}）。
  \item \textbf{对象池固定，查找是线性扫描。} 128 个 dentry 和 128 个 inode 对现在的用户程序足够。FAT32 自己只有 16 个 \code{fat\_inode} 槽位，打开的文件和它路径上的子目录各占一个。
  \item \textbf{大小存了三份。} \code{inode->i\_size}、\code{i\_backend\_node.size} 和 FAT32 的\code{fat\_inode.size} 由 write 路径分别更新，以后有了真正的 FAT32 \code{inode\_operations} 可以合成一份。
\end{enumerate}

\section{下一步可以怎样演进}

\begin{txtcode}
现在                                       目标
backend_iops（所有 fs 共用）               minifs_iops / fat32_iops / devfs_iops
  lookup(child->d_path)                      lookup(dir->i_ino, child->d_name)
vfs_node 句柄                              inode->i_private 直接指向后端 inode
find_mount() 字符串前缀                     逐级走 dentry，遇挂载点切换到 mnt_root
dentry 引用归零即释放                       LRU + 负 dentry
d_path[128] 作为键                          (d_parent, d_name) 哈希
\end{txtcode}

按这个顺序做，每一步都能单独上板测：先修第~\ref{sec:vfs-issues} 节第 1 条让系统能启动，再给minifs 写逐级 lookup（它只读，最简单），然后是 FAT32，最后去掉\code{vfs\_backend\_ops}。
